\section{Discussion}
\label{sec:discussion}

\subsection{Limitations}
The load-bearing limitation is the numerator. The copy time of \eqref{eq:copy} is computed
from a lunar-regolith seed because no per-module mass and energy breakdown yet exists for a
concrete interstellar-probe design; the most developed such design reports its architecture
but not a closed bill of materials at the fidelity this calculation
needs~\cite{borgue-hein-2020}. We therefore do not present 582~days as a probe's copy time.
The paper's whole defense against this gap is the margin of Fig.~\ref{fig:margin}: the
conclusion holds for any copy time within four orders of magnitude of the proxy, which
comfortably spans the range a heavier, less-closed, or more power-starved probe could plausibly
occupy. Closing the gap with a sourced probe bill of materials would replace the proxy with a
measurement, but on the present evidence it cannot change the verdict.

The dwell is evaluated at 1~AU around a Sun-like star, matching the uniform field the
settlement-front model assumes but not the real spread of stellar luminosities and material
availability. Because stellar variation enters only through the numerator, its effect is
already bounded by the same margin; a fainter or resource-poorer star lengthens the copy time,
and it would have to do so by four orders of magnitude to matter. We deliberately do not add a
stellar-population sub-model: it would be a separate slice with its own validation, not a
parameter of this one, and the margin makes it unnecessary for the present claim. Likewise we
inherit the front model's idealizations, chiefly a frozen stellar field, but our claim is a
ratio in which the dwell enters the same way whether or not those idealizations hold, so it is
insulated from them in a way an absolute fill time would not be.

Finally, the direct tax for max-boost slingshots is unresolved at the field size we can
afford (Fig.~\ref{fig:tax}); we report it as positive but within seed noise rather than as a
number. That honesty is the point of the ensemble. The powered tax is smaller still and is
knowable only analytically, through \eqref{eq:frac}, not by brute force.

\subsection{Implications}
The substantive implication is a division of labour between scales. Manufacturing cadence is
the binding constraint on a local, densely packed fleet, where it sets the doubling time
outright, and is a negligible tax on a sparse interstellar front, where propulsion policy and
raw transit time dominate the timeline by six orders of magnitude. A programme that worries
about how fast its seeds can build is asking the right question for the solar-system phase and
the wrong one for the galactic phase, where the same effort is better spent on propulsion. The
result also illustrates a modelling stance: by keeping the three scales as separate, honestly
sourced models joined at a clean seam, rather than one fused simulation, we can derive a shared
quantity in the place it is physically fixed and route it to the place that needs it, and read
off which constraint binds where. The seam is what makes the cross-scale statement possible;
fusing the scales would have buried it.
